\documentclass[aip,jcp,reprint,noshowkeys,superscriptaddress]{revtex4-1}
\usepackage{graphicx,dcolumn,bm,xcolor,microtype,multirow,amscd,amsmath,amssymb,amsfonts,physics,longtable,wrapfig}
\usepackage[version=4]{mhchem}
\allowdisplaybreaks[1]

\usepackage[utf8]{inputenc}
\usepackage[T1]{fontenc}
\usepackage{txfonts}

\usepackage[
    colorlinks,
    linkcolor={red!50!black},
    citecolor={red!70!black},
    urlcolor={red!80!black}
	]{hyperref}
\urlstyle{same}

\newcommand{\ie}{\textit{i.e.}}
\newcommand{\eg}{\textit{e.g.}}
\newcommand{\alert}[1]{\textcolor{red}{#1}}
\usepackage[normalem]{ulem}
\newcommand{\titou}[1]{\textcolor{red}{#1}}
\newcommand{\trashPFL}[1]{\textcolor{\red}{\sout{#1}}}
\newcommand{\PFL}[1]{\titou{(\underline{\bf PFL}: #1)}}

\newcommand{\cre}[1]{\Hat{a}_{#1}^{\dag}}
\newcommand{\ani}[1]{\Hat{a}_{#1}}
\newcommand{\ca}[2]{\Hat{a}^{#1}_{#2}}

\newcommand{\mc}{\multicolumn}
\newcommand{\fnm}{\footnotemark}
\newcommand{\fnt}{\footnotetext}
\newcommand{\tabc}[1]{\multicolumn{1}{c}{#1}}
\newcommand{\QP}{\textsc{quantum package}}
\newcommand{\T}[1]{#1^{\intercal}}

% coordinates
\newcommand{\br}{\boldsymbol{r}}
\newcommand{\bx}{\boldsymbol{x}}
\newcommand{\bI}{\boldsymbol{1}}
\newcommand{\cK}{\mathcal{K}}
\newcommand{\bG}{\boldsymbol{G}}
\newcommand{\bH}{\boldsymbol{H}}
\newcommand{\bM}{\boldsymbol{M}}
\newcommand{\bE}{\boldsymbol{E}}
\newcommand{\bSigma}{\boldsymbol{\Sigma}}
\newcommand{\bA}{\boldsymbol{A}}
\newcommand{\bB}{\boldsymbol{B}}
\newcommand{\bC}{\boldsymbol{C}}
\newcommand{\bD}{\boldsymbol{D}}
\newcommand{\bS}{\boldsymbol{S}}
\newcommand{\bN}{\boldsymbol{N}}
\newcommand{\bX}{\boldsymbol{X}}
\newcommand{\bY}{\boldsymbol{Y}}
\newcommand{\bZ}{\boldsymbol{Z}}
\newcommand{\bPi}{\boldsymbol{\Pi}}
\newcommand{\bcalM}{\boldsymbol{\mathcal M}}
\newcommand{\bcalS}{\boldsymbol{\mathcal S}}

% orbital energies
\newcommand{\eps}{\epsilon}
\newcommand{\irr}{\text{irr}}

% shortcuts for greek letters
\newcommand{\si}{\sigma}
\newcommand{\la}{\lambda}
\newcommand{\ii}{\mathrm{i}}
\newcommand{\up}{\uparrow}
\newcommand{\dw}{\downarrow}
\newcommand{\upup}{\up\up}
\newcommand{\updw}{\up\dw}
\newcommand{\dwup}{\dw\up}
\newcommand{\dwdw}{\dw\downarrow}

\newcommand{\h}{\text{1h}}
\newcommand{\p}{\text{1p}}
\newcommand{\hp}{\text{1h+1p}}
\newcommand{\hhp}{\text{2h1p}}
\newcommand{\pph}{\text{2p1h}}
\newcommand{\ph}{\text{ph}}
\newcommand{\eh}{\text{eh}}
\newcommand{\he}{\text{he}}
\newcommand{\ee}{\text{ee}}
\newcommand{\teh}{\overline{\text{eh}}}
\newcommand{\pp}{\text{pp}}
\newcommand{\hh}{\text{hh}}
\newcommand{\xc}{\text{xc}}

% this note
\newcommand{\Lag}{\mathcal{L}}
\newcommand{\sLag}{\bar{\mathcal L}}
\newcommand{\act}{\mathrm{act}}
\newcommand{\scv}{\text{cv}}
\newcommand{\sca}{\text{ca}}
\newcommand{\sav}{\text{av}}
\newcommand{\saa}{\text{aa}}
\newcommand{\svv}{\text{vv}}
\newcommand{\scc}{\text{cc}}

% addresses
\newcommand{\LCPQ}{Laboratoire de Chimie et Physique Quantiques (UMR 5626), Universit\'e de Toulouse, CNRS, UPS, France}

% statement-status markers used throughout the notes
\newcommand{\Exact}{\textbf{Exact identity.}\ }
\newcommand{\Established}{\textbf{Established result.}\ }
\newcommand{\Proposed}{\textbf{Proposed approximation.}\ }

\begin{document}

\title{Extended RPA in the particle--hole, particle--particle and crossed particle--hole channels for an arbitrary vacuum:\\ reduced-density-matrix form and active/inactive block structure}

\begin{abstract}
The three single-channel random-phase problems of the fixed-kernel channel model of
\texttt{Pa2CCD.tex} (direct electron--hole, transverse or crossed-ring electron--hole,
and particle--particle) are written in their extended form, \ie\ as Rowe
equations of motion on an arbitrary correlated vacuum $|\Psi_0\rangle$.  All matrix
elements are exact functionals of the one- and two-particle reduced density
matrices of the vacuum only.  Compact closed forms are given in terms of the
generalized Fock matrix, the energy-weighted density (Lagrangian) matrix, the
antisymmetrized integrals and the 2-RDM or its cumulant.  For a CAS-type vacuum
(closed core $\otimes$ correlated active space $\otimes$ empty virtuals) every
matrix is then resolved into the blocks that carry inactive indices only,
active indices only, and mixed indices.  All formulas have been verified
numerically against brute-force Fock-space double commutators.
\end{abstract}

\maketitle

%%%%%%%%%%%%%%%%%%%%%%%%%
\section{Scope and conventions}
\label{sec:conventions}
%%%%%%%%%%%%%%%%%%%%%%%%%

\texttt{Pa2CCD.tex} reduces the parquet channel equations with fixed bare
kernels to three static random-phase problems, one per correlation channel
(direct electron--hole $\eh$, transverse electron--hole $\teh$, particle--particle
$\pp$), all built on a Hartree--Fock determinant.  These notes give the
corresponding equations for an arbitrary vacuum $|\Psi_0\rangle$ in the
``extended'' form of Rowe's equation-of-motion method
\cite{Rowe1968,ChatterjeePernal2012}: the pair operators of each channel act on
the correlated vacuum, the double commutator with $\hat H$ replaces the bare
kernel, and the commutator metric replaces the Hartree--Fock pair norm.  Every
matrix element is then an exact functional of the one- and two-particle reduced
density matrices (RDMs) of $|\Psi_0\rangle$.  The second part
(Sec.~\ref{sec:blocks}) resolves these matrices into inactive, active and mixed
blocks for the CAS-type vacuum used throughout \texttt{mr-hedin.tex}.

\paragraph*{Hamiltonian and integrals.}
All quantities are real (real spin orbitals, real vacuum).  With
$\langle pq|rs\rangle=(pr|qs)=v_{pr,qs}$ the notation of \texttt{mr-hedin.tex},
\begin{align}
  \hat H&=\sum_{pq}h_{pq}\,\cre{p}\ani{q}
  +\frac12\sum_{pqrs}\langle pq|rs\rangle\,\cre{p}\cre{q}\ani{s}\ani{r},
  \nonumber\\
  \langle pq||rs\rangle&=\langle pq|rs\rangle-\langle pq|sr\rangle
  =v_{pr,qs}-v_{ps,qr}.
  \label{eq:H}
\end{align}
General spin orbitals are $p,q,r,s$ and general summation indices $t,u$;
in Sec.~\ref{sec:blocks} $i,j,k,l$ denote core, $a,b,c,d$ virtual and
$x,y,z,w$ active spin orbitals, and the summation indices $u,v$ run over
active spin orbitals only.

\paragraph*{Reduced density matrices.}
\begin{align}
  \gamma_{pq}&=\langle\Psi_0|\cre{p}\ani{q}|\Psi_0\rangle,
  \qquad
  \Gamma_{pq,rs}=\langle\Psi_0|\cre{p}\cre{q}\ani{s}\ani{r}|\Psi_0\rangle,
  \label{eq:rdm}\\
  \lambda_{pq,rs}&=\Gamma_{pq,rs}-\gamma_{pr}\gamma_{qs}+\gamma_{ps}\gamma_{qr},
  \label{eq:cumulant}
\end{align}
so that $E_0=\sum_{pq}h_{pq}\gamma_{pq}+\tfrac12\sum v_{pr,qs}\Gamma_{pq,rs}$,
$\Gamma_{pq,rs}=-\Gamma_{qp,rs}=-\Gamma_{pq,sr}=\Gamma_{rs,pq}$ and
$\lambda$ has the same symmetries.  $\Gamma$ coincides with $D^{(2)}$ of
\texttt{numerics/equations.tex} [its Eq.~(rdm-def)] and with
$\Gamma_{pq,sr}$ in the notation of the EKT section of \texttt{mr-hedin.tex}.
Nothing is assumed about $|\Psi_0\rangle$ except that it is a normalized
$N$-electron state; in particular it need not be an eigenstate of $\hat H$.

\paragraph*{Two one-body matrices.}
All closed forms below are expressed with the generalized Fock matrix $f$ and
the energy-weighted density (Lagrangian) matrix $\Lag$,
\begin{align}
  f_{pq}&=h_{pq}+\sum_{tu}\gamma_{tu}\,\langle pt||qu\rangle ,
  \label{eq:genfock}\\
  \Lag_{pq}&=\sum_{t}h_{pt}\gamma_{tq}+\frac12\sum_{tuv}\langle pu||tv\rangle\,\Gamma_{qu,tv}
  \nonumber\\
  &=\sum_t f_{pt}\gamma_{tq}+\frac12\sum_{tuv}\langle pu||tv\rangle\,\lambda_{qu,tv},
  \label{eq:lagrangian}\\
  \sLag_{pq}&=\tfrac12\bigl(\Lag_{pq}+\Lag_{qp}\bigr).
  \label{eq:symlag}
\end{align}
$f$ is symmetric; $\Lag$ is symmetric only if $|\Psi_0\rangle$ is stationary
with respect to orbital rotations, because
\begin{equation}
  \langle\Psi_0|[\hat H,\cre{p}\ani{q}]|\Psi_0\rangle=\Lag_{pq}-\Lag_{qp}
  \label{eq:brillouin}
\end{equation}
(generalized Brillouin condition).  $\mathrm{Tr}\,\Lag=E_1+2E_2$.

\paragraph*{Status markers.}
As in \texttt{mr-hedin.tex}: \Exact exact identities; \Established results
available in the literature or in the other notes; \Proposed definitions made
here.

%%%%%%%%%%%%%%%%%%%%%%%%%
\section{Equation of motion in an operator manifold}
\label{sec:eom}
%%%%%%%%%%%%%%%%%%%%%%%%%

\Established For a set of (even) operators $\{\hat Q_\mu^\dagger\}$ define
\begin{align}
  \mathcal M_{\mu\nu}&=\langle\Psi_0|[\hat Q_\mu,[\hat H,\hat Q_\nu^\dagger]]|\Psi_0\rangle ,
  \qquad
  \mathcal S_{\mu\nu}=\langle\Psi_0|[\hat Q_\mu,\hat Q_\nu^\dagger]|\Psi_0\rangle ,
  \label{eq:MS-def}\\
  \bcalM\,c_\nu&=\omega_\nu\,\bcalS\,c_\nu ,\qquad
  c_\nu^{T}\bcalS c_{\nu'}=\pm\delta_{\nu\nu'} .
  \label{eq:eom}
\end{align}
If $|\Psi_0\rangle$ is the exact ground state and the manifold contains the
exact excitation operators $|n\rangle\langle 0|$, the solutions of
Eq.~\eqref{eq:eom} are the exact excitation energies and the
positive-norm (negative-norm) eigenvectors give the exact transition amplitudes
$\langle n|\hat Q^\dagger_\nu|0\rangle$ (de-excitation amplitudes); for a
truncated manifold Eq.~\eqref{eq:eom} is the inner projection of the
superoperator resolvent onto the manifold, and the corresponding pair
propagator is
\begin{equation}
  \bPi(\Omega)=\bcalS\,(\Omega\bcalS-\bcalM)^{-1}\bcalS
  =\sum_{\nu}\frac{(\bcalS c_\nu)(\bcalS c_\nu)^{T}}{c_\nu^{T}\bcalS c_\nu\,(\Omega-\omega_\nu)} .
  \label{eq:propagator}
\end{equation}
Equation~\eqref{eq:propagator} is the extended-vacuum counterpart of the static
channel propagators
$\boldsymbol L^{\alpha}(\Omega)=\ii\sigma_\alpha(\Omega\boldsymbol\eta^{\alpha}-\boldsymbol{\mathcal H}^{\alpha})^{-1}$
of \texttt{Pa2CCD.tex}: $\bcalS$ plays the role of the sector metric
$\boldsymbol\eta^\alpha$ (it is no longer $\pm\bI$) and $\bcalM$ that of
$\boldsymbol{\mathcal H}^\alpha=\boldsymbol{\mathcal H}^\alpha_0+\boldsymbol V^\alpha$.

\Exact For a vacuum that is not an eigenstate of $\hat H$, the double
commutator \eqref{eq:MS-def} is not symmetric; by the Jacobi identity
$\mathcal M_{\mu\nu}-\mathcal M_{\nu\mu}=\langle[\hat H,[\hat Q_\mu,\hat Q^\dagger_\nu]]\rangle$,
which for the pair manifolds below reduces to Kronecker deltas times
$\Lag-\Lag^{T}$ [see Eq.~\eqref{eq:ph-asym}].  Throughout, $\bcalM$ denotes the
symmetric part $\tfrac12(\bcalM+\bcalM^{T})$, which is what the implementation
in \texttt{numerics/mrgw\_hermitian/ekt\_erpa.py} uses; in all closed forms this
amounts to replacing $\Lag$ by $\sLag$.

%%%%%%%%%%%%%%%%%%%%%%%%%
\section{Particle--hole channel (\texorpdfstring{$\eh$}{eh}, ring)}
\label{sec:ph}
%%%%%%%%%%%%%%%%%%%%%%%%%

\Established The manifold consists of all single replacements
$\hat E_{rs}=\cre{r}\ani{s}$, $r\ne s$, on the correlated vacuum; the
composite index $(pq)$ labels the bra operator $\hat E_{pq}$, whose adjoint
$\hat E_{qp}$ appears in the commutators, \ie\
\begin{equation}
  \mathcal M_{(pq),(rs)}=\langle[\hat E_{qp},[\hat H,\hat E_{rs}]]\rangle ,
  \qquad
  \mathcal S_{(pq),(rs)}=\langle[\hat E_{qp},\hat E_{rs}]\rangle .
  \label{eq:ph-def}
\end{equation}
This is the ERPA of Chatterjee and Pernal \cite{ChatterjeePernal2012} and the
active-space ERPA of \texttt{mr-hedin.tex} [its Eq.~(erpa-eom)], here written
in the full orbital space.

\Exact The metric and the (symmetrized) double commutator are
\begin{widetext}
\begin{align}
  \mathcal S_{(pq),(rs)}&=\delta_{pr}\gamma_{qs}-\delta_{qs}\gamma_{pr},
  \label{eq:ph-S}\\
  \mathcal M_{(pq),(rs)}&=
  h_{pr}\gamma_{qs}+h_{qs}\gamma_{pr}-\delta_{pr}\sLag_{qs}-\delta_{qs}\sLag_{pr}
  +\sum_{tu}\Bigl[\langle pt||ur\rangle\,\Gamma_{qt,us}+\langle qt||us\rangle\,\Gamma_{pt,ur}\Bigr]
  -\frac12\sum_{tu}\Bigl[\langle ps||tu\rangle\,\Gamma_{qr,tu}+\langle qr||tu\rangle\,\Gamma_{ps,tu}\Bigr]
  \label{eq:ph-M}\\
  &=\gamma_{qs}f_{pr}+\gamma_{pr}f_{qs}-\delta_{pr}\sLag_{qs}-\delta_{qs}\sLag_{pr}
  +\sum_{tu}\Bigl[\langle pt||ur\rangle\gamma_{qu}\gamma_{ts}+\langle qt||us\rangle\gamma_{pu}\gamma_{tr}
  -\langle ps||tu\rangle\gamma_{qt}\gamma_{ru}-\langle qr||tu\rangle\gamma_{pt}\gamma_{su}\Bigr]
  \nonumber\\
  &\quad
  +\sum_{tu}\Bigl[\langle pt||ur\rangle\lambda_{qt,us}+\langle qt||us\rangle\lambda_{pt,ur}\Bigr]
  -\frac12\sum_{tu}\Bigl[\langle ps||tu\rangle\lambda_{qr,tu}+\langle qr||tu\rangle\lambda_{ps,tu}\Bigr].
  \label{eq:ph-M-cum}
\end{align}
\end{widetext}
Equation~\eqref{eq:ph-M} is the 20-term list of \texttt{numerics/rdm\_terms.tex}
(generated by \texttt{wick.py}) after collecting the direct and exchange
integrals; the unsymmetrized double commutator is obtained by replacing
$\sLag_{qs}\to\Lag_{sq}$ and $\sLag_{pr}\to\Lag_{rp}$, so that
\begin{equation}
  \mathcal M_{(pq),(rs)}-\mathcal M_{(rs),(pq)}
  =-\delta_{pr}(\Lag_{sq}-\Lag_{qs})-\delta_{qs}(\Lag_{rp}-\Lag_{pr})
  \label{eq:ph-asym}
\end{equation}
before symmetrization.  The cumulant form \eqref{eq:ph-M-cum} makes the
mean-field content explicit: for a single determinant ($\lambda=0$,
$\gamma^2=\gamma$) it collapses to the usual RPA matrices, see below.
Two further exact properties, both used in the next paragraph, are
$\mathcal M_{(qp),(sr)}=\mathcal M_{(pq),(rs)}$ (adjoint of the double
commutator with real quantities) and the block-diagonality of $\bcalS$ in the
natural-orbital basis.

\Established \emph{RPA form.}  In the natural spin-orbital basis
$\gamma=\mathrm{diag}(n_p)$ the metric is diagonal,
$\mathcal S_{(pq),(rs)}=\delta_{pr}\delta_{qs}(n_q-n_p)$.  Let
$\mathcal E=\{(pq):n_q>n_p\}$ be the excitation pairs (particle-like index
first), $\bar{\mathcal E}=\{(qp)\}$ their de-excitation partners, and
$N_{(pq)}=n_q-n_p>0$.  With
\begin{equation}
  A_{(pq),(rs)}=\mathcal M_{(pq),(rs)},\;\;
  B_{(pq),(rs)}=-\mathcal M_{(pq),(sr)},\;\; (pq),(rs)\in\mathcal E,
  \label{eq:ph-AB}
\end{equation}
both symmetric, Eq.~\eqref{eq:eom} with
$\hat Q^\dagger_\nu=\sum_{\mathcal E}\bigl(X_{\nu,pq}\hat E_{pq}-Y_{\nu,pq}\hat E_{qp}\bigr)$
becomes
\begin{equation}
  \begin{pmatrix}\bA&\bB\\ \bB&\bA\end{pmatrix}
  \begin{pmatrix}\bX_\nu\\ \bY_\nu\end{pmatrix}
  =\omega_\nu
  \begin{pmatrix}\bN&0\\0&-\bN\end{pmatrix}
  \begin{pmatrix}\bX_\nu\\ \bY_\nu\end{pmatrix},
  \label{eq:ph-rpa}
\end{equation}
\ie\ the standard RPA structure with the metric $\bN$; the transition
densities are $\rho^{\nu}_{pq}=N_{(pq)}X_{\nu,pq}$ and
$\rho^{\nu}_{qp}=-N_{(pq)}Y_{\nu,pq}$ [Eq.~(erpa-rho) of
\texttt{mr-hedin.tex}, whose $\tilde B$ equals $-\bB$ because its
de-excitation coefficients are $-\bY$].  Pairs with $n_p=n_q$ have zero norm and
are removed from the manifold.  For a Hartree--Fock vacuum $\mathcal E=\{(ai)\}$
and Eqs.~\eqref{eq:ph-M-cum} and \eqref{eq:ph-AB} give
$A_{(ai),(bj)}=\delta_{ij}f_{ab}-\delta_{ab}f_{ij}+\langle aj||ib\rangle$,
$B_{(ai),(bj)}=\langle ab||ij\rangle$, the matrices of the direct
electron--hole channel of \texttt{Pa2CCD.tex} [its $\bA^{\eh}$, $\bB^{\eh}$].

%%%%%%%%%%%%%%%%%%%%%%%%%
\section{Particle--particle channel (\texorpdfstring{$\pp$}{pp}, ladder)}
\label{sec:pp}
%%%%%%%%%%%%%%%%%%%%%%%%%

\Established The manifold consists of all pair-creation operators
$\cre{r}\cre{s}$, $r<s$, on the correlated vacuum (no distinction between
particle and hole pairs is made at this stage):
\begin{equation}
  \mathcal M^{\pp}_{(pq),(rs)}=\langle[\ani{q}\ani{p},[\hat H,\cre{r}\cre{s}]]\rangle,
  \qquad
  \mathcal S^{\pp}_{(pq),(rs)}=\langle[\ani{q}\ani{p},\cre{r}\cre{s}]\rangle .
  \label{eq:pp-def}
\end{equation}
Both are antisymmetric under $p\leftrightarrow q$ and under
$r\leftrightarrow s$.  With $P(pq)X_{pq}=X_{pq}-X_{qp}$,

\Exact
\begin{widetext}
\begin{align}
  \mathcal S^{\pp}_{(pq),(rs)}&=\delta_{pr}\delta_{qs}-\delta_{ps}\delta_{qr}
  -\delta_{pr}\gamma_{qs}+\delta_{ps}\gamma_{qr}+\delta_{qr}\gamma_{ps}-\delta_{qs}\gamma_{pr},
  \label{eq:pp-S}\\
  \mathcal M^{\pp}_{(pq),(rs)}&=
  P(pq)P(rs)\Bigl[\delta_{qs}\bigl(f_{pr}-\sLag_{pr}\bigr)-h_{pr}\gamma_{qs}
  +\sum_{tu}\langle pt||ru\rangle\,\Gamma_{ts,qu}\Bigr]
  \nonumber\\
  &\quad+\langle pq||rs\rangle
  -\sum_t\Bigl[\gamma_{pt}\langle tq||rs\rangle+\gamma_{qt}\langle pt||rs\rangle
  +\gamma_{rt}\langle pq||ts\rangle+\gamma_{st}\langle pq||rt\rangle\Bigr]
  \label{eq:pp-M}\\
  &=P(pq)P(rs)\Bigl[\delta_{qs}\bigl(f_{pr}-\sLag_{pr}\bigr)-\gamma_{qs}f_{pr}
  +\sum_{tu}\langle pt||ru\rangle\gamma_{qt}\gamma_{su}
  +\sum_{tu}\langle pt||ru\rangle\lambda_{ts,qu}\Bigr]
  \nonumber\\
  &\quad+\langle pq||rs\rangle
  -\sum_t\Bigl[\gamma_{pt}\langle tq||rs\rangle+\gamma_{qt}\langle pt||rs\rangle
  +\gamma_{rt}\langle pq||ts\rangle+\gamma_{st}\langle pq||rt\rangle\Bigr].
  \label{eq:pp-M-cum}
\end{align}
\end{widetext}
Equation~\eqref{eq:pp-M} collects the 72 terms produced by \texttt{wick.py} for
the pair manifold; the unsymmetrized double commutator has $\Lag_{rp}$ in place
of $\sLag_{pr}$.  Note that the 2-RDM enters only through particle--hole-type
contractions $\sum_{tu}\langle pt||ru\rangle\Gamma_{ts,qu}$ and through $\Lag$;
the ladder-type integral $\langle pq||rs\rangle$ is dressed by the 1-RDM only
(one factor $-\gamma$ on each of the four legs).

\Established \emph{RPA form.}  In the natural spin-orbital basis, for $p<q$
and $r<s$, $\mathcal S^{\pp}_{(pq),(rs)}=\delta_{pr}\delta_{qs}(1-n_p-n_q)$.
Pairs with $1-n_p-n_q>0$ form the particle-like sector $\mathcal P$, pairs with
$1-n_p-n_q<0$ the hole-like sector $\mathcal H$, and pairs with $n_p+n_q=1$
(\eg\ one fully occupied and one empty orbital) have zero norm and are removed.
With $N^{\pp}_{(pq)}=|1-n_p-n_q|$ and
\begin{equation}
  \bA^{\pp}=\bcalM^{\pp}_{\mathcal P\mathcal P},\qquad
  \bC^{\pp}=\bcalM^{\pp}_{\mathcal H\mathcal H},\qquad
  \bB^{\pp}=-\bcalM^{\pp}_{\mathcal P\mathcal H},
  \label{eq:pp-ABC}
\end{equation}
Eq.~\eqref{eq:eom} with
$\hat Q^\dagger_\nu=\sum_{\mathcal P}X_{\nu,pq}\cre{p}\cre{q}-\sum_{\mathcal H}Y_{\nu,pq}\cre{p}\cre{q}$
reads
\begin{equation}
  \begin{pmatrix}\bA^{\pp}&\bB^{\pp}\\ \bB^{\pp T}&\bC^{\pp}\end{pmatrix}
  \begin{pmatrix}\bX_\nu\\ \bY_\nu\end{pmatrix}
  =\omega_\nu
  \begin{pmatrix}\bN^{\pp}_{\mathcal P}&0\\0&-\bN^{\pp}_{\mathcal H}\end{pmatrix}
  \begin{pmatrix}\bX_\nu\\ \bY_\nu\end{pmatrix}.
  \label{eq:pp-rpa}
\end{equation}
Positive-norm solutions are the $N{+}2$ states, $\omega_\nu=E^{N+2}_\nu-E_0$;
negative-norm solutions are the $N{-}2$ states with
$\omega_\nu=-(E^{N-2}_\nu-E_0)$ \cite{SchuckBook,Scuseria2013}.  The sign in
the definition of $\bB^{\pp}$ compensates the sign of the commutator metric in
$\mathcal H$ and makes the Hartree--Fock limit coincide with the
particle--particle channel of \texttt{Pa2CCD.tex}: $\mathcal P=\{(ab)\}$,
$\mathcal H=\{(ij)\}$ and Eq.~\eqref{eq:pp-M-cum} gives, in canonical orbitals,
$A^{\pp}_{ab,cd}=(\eps_a+\eps_b)\delta^{\rm A}_{ab,cd}+\langle ab||cd\rangle$,
$C^{\pp}_{ij,kl}=-(\eps_i+\eps_j)\delta^{\rm A}_{ij,kl}+\langle ij||kl\rangle$,
$B^{\pp}_{ab,ij}=\langle ab||ij\rangle$ [its $\bC^{\pp}$, $\bD^{\pp}$, $\bB^{\pp}$
with ordered pairs].

%%%%%%%%%%%%%%%%%%%%%%%%%
\section{Crossed-ring particle--hole channel (\texorpdfstring{$\teh$}{eh-bar})}
\label{sec:teh}
%%%%%%%%%%%%%%%%%%%%%%%%%

\Established The transverse channel is the electron--hole channel with
crossed external labels: the composite pair joins the hole label of one
external vertex with the particle label of the other, and the fermionic
crossing sign
$\hat E_{bi}\hat E_{aj}|\Phi_0\rangle=-\hat E_{ai}\hat E_{bj}|\Phi_0\rangle$
[Eq.~(crossed\_operator\_sign) of \texttt{Pa2CCD.tex}] is carried by the
forward--backward blocks.  In the Hartree--Fock static channel reduction of
\texttt{Pa2CCD.tex} the generator in the crossed composite ordering
$(ib),(ja)$ is
\begin{align}
  \bH^{\teh}&=\begin{pmatrix}\bA^{\teh}&-\bB^{\teh}\\ \bB^{\teh\dagger}&-\bA^{\teh\dagger}\end{pmatrix},
  \nonumber\\
  \bA^{\teh}_{ib,kc}&=(\eps_b-\eps_i)\delta_{ik}\delta_{bc}+\langle bk||ic\rangle,
  \qquad
  \bB^{\teh}_{ib,ja}=\langle ba||ij\rangle .
  \label{eq:H-teh-HF}
\end{align}
Read as arrays over (hole, particle) pairs, $\bA^{\teh}$ and $\bB^{\teh}$ are
identical to the direct-channel matrices $\bA^{\eh}_{ia,jb}=(\eps_a-\eps_i)\delta_{ij}\delta_{ab}+\langle aj||ib\rangle$
and $\bB^{\eh}_{ia,jb}=\langle ab||ij\rangle$; the crossing is carried entirely
by the sign of the forward--backward block and by the crossed attachment of
the composite indices to the external legs (the map $J_{\teh}$ of
\texttt{Pa2CCD.tex}).  With $\bZ=\mathrm{diag}(\bI,-\bI)$ on the
forward/backward sectors,
$\bH^{\teh}=\bZ\bH^{\eh}\bZ$ and $\boldsymbol L^{\teh}(\Omega)=\bZ\boldsymbol L^{\eh}(\Omega)\bZ$:
the crossed channel has the same poles as the direct one, the eigenvectors
$(\bX_\nu,-\bY_\nu)$, and its static propagator differs from the direct one by
the sign of the forward--backward blocks.  The channels differ once they are
attached to external legs; in the common-amplitude construction this is what
turns the ring terms of the $\eh$ channel into the crossed-ring terms of the
$\teh$ channel [Eqs.~(ring\_linear\_audit) and (ring\_quadratic\_audit) of
\texttt{Pa2CCD.tex}].

\Proposed \emph{Extended form.}  The same construction defines the crossed-ring
channel on an arbitrary vacuum.  Let $\bZ=\bcalS\,(\bcalS^{2})^{-1/2}$ be the
sign of the particle--hole metric on its nonzero-norm subspace
($\bZ=\mathrm{diag}(\bI,-\bI)$ on $\mathcal E\oplus\bar{\mathcal E}$ in the
natural-orbital basis).  The crossed-ring ERPA is
\begin{align}
  \bcalM^{\teh}&=\bZ\,\bcalM\,\bZ,\qquad \bcalS^{\teh}=\bZ\bcalS\bZ=\bcalS,
  \nonumber\\
  \begin{pmatrix}\bA&-\bB\\ -\bB&\bA\end{pmatrix}
  \begin{pmatrix}\bX_\nu\\ \bY_\nu\end{pmatrix}
  &=\omega_\nu
  \begin{pmatrix}\bN&0\\0&-\bN\end{pmatrix}
  \begin{pmatrix}\bX_\nu\\ \bY_\nu\end{pmatrix},
  \label{eq:teh-erpa}
\end{align}
with $\bA$, $\bB$, $\bN$ of Eqs.~\eqref{eq:ph-AB}--\eqref{eq:ph-rpa}, and its
propagator is $\bPi^{\teh}(\Omega)=\bZ\bPi(\Omega)\bZ$.  No matrix elements
beyond those of Sec.~\ref{sec:ph} are needed: the 1-/2-RDM content of the
crossed channel is that of Eqs.~\eqref{eq:ph-S}--\eqref{eq:ph-M-cum}, with the
sign of $\bB$ reversed.  Equation~\eqref{eq:teh-erpa} reduces to
Eq.~\eqref{eq:H-teh-HF} for a Hartree--Fock vacuum and is the unique
extension for which the exact crossing relation
$L(12;1'2')=-L(12;2'1')$ of \texttt{parquet.tex} acts on the static
channel propagator in the same way as in the Hartree--Fock case, namely on
the forward/backward structure only.  (With the antisymmetrized kernel both
the direct and the exchange integrals are present in every channel; ``ring''
and ``crossed ring'' refer to the pairing of the external legs, not to a
direct/exchange split of the kernel.)

%%%%%%%%%%%%%%%%%%%%%%%%%
\section{Block structure for a CAS-type vacuum}
\label{sec:blocks}
%%%%%%%%%%%%%%%%%%%%%%%%%

\subsection{Structure of the vacuum and of \texorpdfstring{$f$, $\Lag$}{f, L}}
\label{sec:cas-structure}

\Exact Let $|\Psi_0\rangle=|\Phi_{\rm core}\rangle\otimes|\Xi\rangle$ with all
core spin orbitals $i,j,\ldots$ occupied, all virtual spin orbitals
$a,b,\ldots$ empty, and $|\Xi\rangle$ an \emph{arbitrary} $N_{\rm act}$-electron
state in the active spin orbitals $x,y,\ldots$ (CASCI or CASSCF, but no
stationarity is assumed).  Then
\begin{align}
  \gamma_{ij}&=\delta_{ij},\quad \gamma_{ab}=0,\quad \gamma_{xy}=\gamma^{\act}_{xy},\quad
  \gamma_{Px}=\gamma_{ia}=0,
  \nonumber\\
  \lambda_{pq,rs}&\ne0\quad\text{only for }p,q,r,s\in\text{active},
  \label{eq:cas-rdm}
\end{align}
and $\Gamma$ follows from Eq.~\eqref{eq:cumulant}.  Inserting
Eq.~\eqref{eq:cas-rdm} into Eqs.~\eqref{eq:genfock}--\eqref{eq:symlag},
with $u,v,z$ active,
\begin{align}
  f_{pq}&=h_{pq}+\sum_{k}\langle pk||qk\rangle+\sum_{uv}\gamma_{uv}\langle pu||qv\rangle ,
  \label{eq:cas-fock}\\
  \Lag_{pj}&=f_{pj},\qquad \Lag_{pb}=0,\qquad
  \Lag_{py}=\sum_{u}f_{pu}\gamma_{uy}+\Lambda_{py},
  \nonumber\\
  \Lambda_{py}&=\frac12\sum_{uvz}\langle pu||vz\rangle\,\lambda_{yu,vz},
  \label{eq:cas-lag}
\end{align}
so that
\begin{align}
  \sLag_{ij}&=f_{ij},\qquad \sLag_{ab}=0,\qquad \sLag_{ia}=\tfrac12 f_{ia},
  \nonumber\\
  \sLag_{ix}&=\tfrac12\Bigl[f_{xi}+\sum_u f_{iu}\gamma_{ux}+\Lambda_{ix}\Bigr],
  \nonumber\\
  \sLag_{ax}&=\tfrac12\Bigl[\sum_u f_{au}\gamma_{ux}+\Lambda_{ax}\Bigr],
  \nonumber\\
  \sLag_{xy}&=\tfrac12\Bigl[\sum_u\bigl(f_{xu}\gamma_{uy}+f_{yu}\gamma_{ux}\bigr)+\Lambda_{xy}+\Lambda_{yx}\Bigr].
  \label{eq:cas-symlag}
\end{align}
$f$ is the generalized (Dyall) Fock matrix that also defines the inactive
orbital energies $\kappa_P=\eps_P$ in \texttt{mr-hedin.tex}, and its active
block is $f_{xy}=h^{\rm eff}_{xy}+\sum_{uv}\gamma_{uv}\langle xu||yv\rangle$ with
$h^{\rm eff}_{xy}=h_{xy}+\sum_k\langle xk||yk\rangle$ the one-body part of the
active Dyall Hamiltonian; $\Lag_{xy}$ is therefore the Lagrangian of the
active-space problem.  For a CASSCF vacuum $\Lag$ is symmetric
[Eq.~\eqref{eq:brillouin}]: $f_{ia}=0$, $f_{xi}=\sum_u f_{iu}\gamma_{ux}+\Lambda_{ix}$,
$\sum_u f_{au}\gamma_{ux}+\Lambda_{ax}=0$, and the bars in
Eq.~\eqref{eq:cas-symlag} can be dropped; for CASCI orbitals the
symmetrization is part of the definition of $\bcalM$.

In the block formulas below every summation index $u,v$ runs over the
\emph{active} spin orbitals only: the core contractions have been absorbed
into $f$ and the Kronecker deltas, and the virtual ones vanish because
$\gamma_{ab}=0$.  The blocks are exact consequences of
Eqs.~\eqref{eq:ph-M-cum}, \eqref{eq:pp-M-cum} and \eqref{eq:cas-rdm}; each of
them was generated by substituting the index types into the general formulas
and verified numerically (Sec.~\ref{sec:verification}).  Every block depends
on the vacuum only through $\gamma^{\act}$, $\lambda^{\act}$ (equivalently
$\Gamma^{\act}$), $f$ and $\sLag$.

\subsection{Particle--hole channel}
\label{sec:ph-blocks}

\Exact The pairs with nonzero norm fall into four sectors, which are the four
channels of the MR-RPA of Ref.~\cite{WangFangLi2026},
\begin{equation}
  \scv:\ (ai),\qquad \sca:\ (xi),\qquad \sav:\ (ax),\qquad \saa:\ (xy),\ x\ne y ,
  \label{eq:ph-sectors}
\end{equation}
with their de-excitation partners $(ia)$, $(ix)$, $(xa)$; the pairs $(ij)$ and
$(ab)$ have zero norm.  The metric is block diagonal in the sectors,
\begin{align}
  \mathcal S_{(ai),(bj)}&=\delta_{ab}\delta_{ij},\qquad
  \mathcal S_{(xi),(yj)}=\delta_{ij}\bigl(\delta_{xy}-\gamma_{xy}\bigr),
  \nonumber\\
  \mathcal S_{(ax),(by)}&=\delta_{ab}\gamma_{xy},\qquad
  \mathcal S_{(xy),(zw)}=\delta_{xz}\gamma_{yw}-\delta_{yw}\gamma_{xz},
  \label{eq:ph-S-blocks}
\end{align}
(in active natural orbitals: $1$, $1-n_x$, $n_x$ and $n_y-n_x$), whereas
$\bcalM$ couples all sectors.  For $\scv$, $\sca$ and $\sav$ the
excitation/de-excitation assignment is fixed, and we list $A$ and $B$ of
Eq.~\eqref{eq:ph-AB}; for the $\saa$ sector it requires the active natural
orbitals, so for all blocks involving $\saa$ the full
$\mathcal M_{(pq),(zw)}$ is given for all ordered pairs $(zw)$, from which
$A$ and $B$ follow by Eq.~\eqref{eq:ph-AB} after the transformation.

\onecolumngrid
\paragraph*{Inactive indices only ($\scv$--$\scv$).}
The RPA matrices of the core$\to$virtual excitations with the generalized Fock
operator; no explicit dependence on the active-space correlation beyond
$\gamma^{\act}$ in $f$:
\begin{align}
  A_{(ai),(bj)} &= \delta_{ij}f_{ab} -\delta_{ab}f_{ij} +\langle aj||ib\rangle
  \label{eq:ph-A-cv-cv}\\
  B_{(ai),(bj)} &= \langle ab||ij\rangle
  \label{eq:ph-B-cv-cv}
\end{align}

\paragraph*{Active indices only ($\saa$--$\saa$).}
Equation~\eqref{eq:ph-M-cum} restricted to active indices with $h\to h^{\rm eff}$,
\ie\ the active-space ERPA of \texttt{mr-hedin.tex} [its Eq.~(erpa-eom)] and of
\texttt{ekt\_erpa.py}:
\begin{align}
  M_{(xy),(zw)} &= -\delta_{yw}\bar{\mathcal L}_{xz} -\delta_{xz}\bar{\mathcal L}_{yw} +f_{xz}\gamma_{yw} \nonumber\\
  &\quad +f_{yw}\gamma_{xz} -\sum_{uv}\langle uw||vy\rangle\gamma_{ux}\gamma_{vz} -\sum_{uv}\langle uv||yz\rangle\gamma_{ux}\gamma_{vw} \nonumber\\
  &\quad -\sum_{uv}\langle uv||xw\rangle\gamma_{uy}\gamma_{vz} -\sum_{uv}\langle uz||vx\rangle\gamma_{uy}\gamma_{vw} -\tfrac{1}{2}\sum_{uv}\langle uv||yz\rangle\lambda_{uv,xw} \nonumber\\
  &\quad -\tfrac{1}{2}\sum_{uv}\langle uv||xw\rangle\lambda_{uv,yz} +\sum_{uv}\langle uy||vw\rangle\lambda_{ux,vz} +\sum_{uv}\langle ux||vz\rangle\lambda_{uy,vw}
  \label{eq:ph-A-aa-aa}
\end{align}

\paragraph*{Mixed blocks.}
Diagonal blocks of the sectors with one active index and all inter-sector
couplings:
\begin{align}
  A_{(xi),(yj)} &= \delta_{ij}f_{xy} -\delta_{xy}f_{ij} -\delta_{ij}\bar{\mathcal L}_{xy} \nonumber\\
  &\quad +f_{ij}\gamma_{xy} +\langle xj||iy\rangle -\sum_{u}\langle uj||iy\rangle\gamma_{ux} \nonumber\\
  &\quad +\sum_{u}\langle ui||xj\rangle\gamma_{uy} -\sum_{uv}\langle uj||vi\rangle\gamma_{ux}\gamma_{vy} +\sum_{uv}\langle ui||vj\rangle\lambda_{ux,vy}
  \label{eq:ph-A-ca-ca}\\
  B_{(xi),(yj)} &= \langle xy||ij\rangle -\sum_{u}\langle uy||ij\rangle\gamma_{ux} +\sum_{u}\langle ux||ij\rangle\gamma_{uy} \nonumber\\
  &\quad +\sum_{uv}\langle uv||ij\rangle\gamma_{ux}\gamma_{vy} +\tfrac{1}{2}\sum_{uv}\langle uv||ij\rangle\lambda_{uv,xy}
  \label{eq:ph-B-ca-ca}\\
  A_{(ax),(by)} &= -\delta_{ab}\bar{\mathcal L}_{xy} +f_{ab}\gamma_{xy} -\sum_{uv}\langle ub||va\rangle\gamma_{ux}\gamma_{vy} \nonumber\\
  &\quad +\sum_{uv}\langle ua||vb\rangle\lambda_{ux,vy}
  \label{eq:ph-A-av-av}\\
  B_{(ax),(by)} &= \sum_{uv}\langle uv||ab\rangle\gamma_{ux}\gamma_{vy} +\tfrac{1}{2}\sum_{uv}\langle uv||ab\rangle\lambda_{uv,xy}
  \label{eq:ph-B-av-av}\\
  A_{(ai),(xj)} &= \delta_{ij}f_{ax} -\delta_{ij}\bar{\mathcal L}_{ax} +\langle aj||ix\rangle \nonumber\\
  &\quad +\sum_{u}\langle ui||aj\rangle\gamma_{ux}
  \label{eq:ph-A-cv-ca}\\
  B_{(ai),(xj)} &= \langle ax||ij\rangle +\sum_{u}\langle ua||ij\rangle\gamma_{ux}
  \label{eq:ph-B-cv-ca}\\
  A_{(ai),(bx)} &= -\delta_{ab}\bar{\mathcal L}_{ix} -\sum_{u}\langle ua||ib\rangle\gamma_{ux}
  \label{eq:ph-A-cv-av}\\
  B_{(ai),(bx)} &= -\sum_{u}\langle ui||ab\rangle\gamma_{ux}
  \label{eq:ph-B-cv-av}\\
  A_{(xi),(ay)} &= -\sum_{u}\langle ux||ia\rangle\gamma_{uy} -\sum_{uv}\langle uv||ia\rangle\gamma_{ux}\gamma_{vy} -\tfrac{1}{2}\sum_{uv}\langle uv||ia\rangle\lambda_{uv,xy}
  \label{eq:ph-A-ca-av}\\
  B_{(xi),(ay)} &= \tfrac{1}{2}\delta_{xy}f_{ia} -f_{ia}\gamma_{xy} -\sum_{u}\langle ui||xa\rangle\gamma_{uy} \nonumber\\
  &\quad +\sum_{uv}\langle ua||vi\rangle\gamma_{ux}\gamma_{vy} -\sum_{uv}\langle ui||va\rangle\lambda_{ux,vy}
  \label{eq:ph-B-ca-av}\\
  M_{(ai),(xy)} &= \sum_{u}\langle ui||ay\rangle\gamma_{ux} -\sum_{u}\langle ua||ix\rangle\gamma_{uy}
  \label{eq:ph-A-cv-aa}\\
  M_{(xi),(yz)} &= -\delta_{xy}\bar{\mathcal L}_{iz} +f_{iz}\gamma_{xy} +\sum_{u}\langle ui||xz\rangle\gamma_{uy} \nonumber\\
  &\quad -\sum_{u}\langle ux||iy\rangle\gamma_{uz} -\sum_{uv}\langle uz||vi\rangle\gamma_{ux}\gamma_{vy} -\sum_{uv}\langle uv||iy\rangle\gamma_{ux}\gamma_{vz} \nonumber\\
  &\quad -\tfrac{1}{2}\sum_{uv}\langle uv||iy\rangle\lambda_{uv,xz} +\sum_{uv}\langle ui||vz\rangle\lambda_{ux,vy}
  \label{eq:ph-A-ca-aa}\\
  M_{(ax),(yz)} &= -\delta_{xz}\bar{\mathcal L}_{ay} +f_{ay}\gamma_{xz} -\sum_{uv}\langle uv||az\rangle\gamma_{ux}\gamma_{vy} \nonumber\\
  &\quad -\sum_{uv}\langle uy||va\rangle\gamma_{ux}\gamma_{vz} -\tfrac{1}{2}\sum_{uv}\langle uv||az\rangle\lambda_{uv,xy} +\sum_{uv}\langle ua||vy\rangle\lambda_{ux,vz}
  \label{eq:ph-A-av-aa}
\end{align}
Remarks.  (i) In $\sca$--$\sca$ and $\sav$--$\sav$ the bare integral is dressed
by the active hole density $\delta_{xy}-\gamma_{xy}$ and by $\gamma$,
respectively, and the cumulant enters through
$\langle ui||vj\rangle\lambda_{ux,vy}$, $\tfrac12\langle uv||ij\rangle\lambda_{uv,xy}$
and $\sLag_{xy}$ (analogously with $a,b$).  The combinations
$\gamma_{ux}\gamma_{vy}+\tfrac12\lambda_{uv,xy}$ contracted with an integral
antisymmetric in $u,v$ equal $\tfrac12\Gamma_{uv,xy}$, so the $B$-type and
$\sca$--$\sav$ blocks can equally be written with the active 2-RDM.
(ii) The couplings between the inactive sector $\scv$ and the active-containing
sectors contain no cumulant: $\scv$--$\sca$ and $\scv$--$\sav$ involve only
$\gamma^{\act}$ and $\sLag$, $\scv$--$\saa$ only $\gamma^{\act}$.  The cumulant
first appears in couplings with at least two active indices in the ket pair or
two active indices across the two pairs.  (iii) For an empty active space only
Eqs.~\eqref{eq:ph-A-cv-cv}--\eqref{eq:ph-B-cv-cv} survive.  For a
determinantal $|\Xi\rangle$ ($\lambda=0$, $\gamma^{\act}$ idempotent) every
block reduces to the corresponding Hartree--Fock RPA block or vanishes together
with its metric.  (iv) The MR-RPA of Ref.~\cite{WangFangLi2026} screens with
the direct residual kernel $v^R$ contracted with the exact channel transition
densities and keeps the exact active response in $\saa$; the blocks above are
instead the single-excitation-manifold double commutator with the
antisymmetrized bare kernel in all four sectors, \ie\ the extended form of the
fixed-kernel model of \texttt{Pa2CCD.tex}.  The two coincide in
$\scv$--$\scv$ (up to direct versus antisymmetrized kernel) and differ elsewhere.

\subsection{Particle--particle channel}
\label{sec:pp-blocks}

\Exact The pairs with nonzero norm fall into five sectors,
\begin{align}
  \svv:\ (ab),\qquad \sav:\ (ax),\qquad \saa:\ (xy),\ x<y,
  \nonumber\\
  \sca:\ (ix),\qquad \scc:\ (ij),
  \label{eq:pp-sectors}
\end{align}
while the core--virtual pairs $(ai)$ have zero norm and drop out.  The metric
is block diagonal,
\begin{align}
  \mathcal S^{\pp}_{(ab),(cd)}&=\delta_{ac}\delta_{bd}-\delta_{ad}\delta_{bc},\qquad
  \mathcal S^{\pp}_{(ax),(by)}=\delta_{ab}\bigl(\delta_{xy}-\gamma_{xy}\bigr),
  \nonumber\\
  \mathcal S^{\pp}_{(ix),(jy)}&=-\delta_{ij}\gamma_{xy},\qquad
  \mathcal S^{\pp}_{(ij),(kl)}=-\bigl(\delta_{ik}\delta_{jl}-\delta_{il}\delta_{jk}\bigr),
  \nonumber\\
  \mathcal S^{\pp}_{(xy),(zw)}&=\delta_{xz}\delta_{yw}-\delta_{xw}\delta_{yz}
  -\delta_{xz}\gamma_{yw}+\delta_{xw}\gamma_{yz}
  \nonumber\\
  &\quad+\delta_{yz}\gamma_{xw}-\delta_{yw}\gamma_{xz},
  \label{eq:pp-S-blocks}
\end{align}
(in active natural orbitals: $1$, $1-n_x$, $1-n_x-n_y$, $-n_x$, $-1$), so that
$\svv$ and $\sav$ belong to the particle-like sector $\mathcal P$, $\scc$ and
$\sca$ to the hole-like sector $\mathcal H$, and $\saa$ splits according to the
sign of $1-n_x-n_y$.  The fifteen blocks of $\bcalM^{\pp}$ are
($P(ab)$ etc.\ are written out):

\paragraph*{Inactive indices only ($\svv$--$\svv$, $\scc$--$\scc$, $\svv$--$\scc$).}
The Hartree--Fock particle--particle, hole--hole and coupling matrices with the
generalized Fock operator [$B^{\pp}_{ab,ij}=-\mathcal M^{\pp}_{(ab),(ij)}=\langle ab||ij\rangle$]:
\begin{align}
  M^{\pp}_{(ab),(cd)} &= \delta_{bd}f_{ac} -\delta_{bc}f_{ad} -\delta_{ad}f_{bc} \nonumber\\
  &\quad +\delta_{ac}f_{bd} +\langle ab||cd\rangle
  \label{eq:pp-vv-vv}\\
  M^{\pp}_{(ij),(kl)} &= -\delta_{jl}f_{ik} +\delta_{jk}f_{il} +\delta_{il}f_{jk} \nonumber\\
  &\quad -\delta_{ik}f_{jl} +\langle ij||kl\rangle
  \label{eq:pp-cc-cc}\\
  M^{\pp}_{(ab),(ij)} &= -\langle ab||ij\rangle
  \label{eq:pp-vv-cc}
\end{align}

\paragraph*{Active indices only ($\saa$--$\saa$).}
Equation~\eqref{eq:pp-M-cum} restricted to active indices with $h\to h^{\rm eff}$,
\ie\ the particle--particle ERPA of the active-space problem:
\begin{align}
  M^{\pp}_{(xy),(zw)} &= \delta_{yw}f_{xz} -\delta_{yz}f_{xw} -\delta_{xw}f_{yz} \nonumber\\
  &\quad +\delta_{xz}f_{yw} -\delta_{yw}\bar{\mathcal L}_{xz} +\delta_{yz}\bar{\mathcal L}_{xw} \nonumber\\
  &\quad +\delta_{xw}\bar{\mathcal L}_{yz} -\delta_{xz}\bar{\mathcal L}_{yw} -f_{xz}\gamma_{yw} \nonumber\\
  &\quad +f_{xw}\gamma_{yz} +f_{yz}\gamma_{xw} -f_{yw}\gamma_{xz} \nonumber\\
  &\quad +\langle xy||zw\rangle -\sum_{u}\langle uy||zw\rangle\gamma_{ux} +\sum_{u}\langle ux||zw\rangle\gamma_{uy} \nonumber\\
  &\quad -\sum_{u}\langle uw||xy\rangle\gamma_{uz} +\sum_{u}\langle uz||xy\rangle\gamma_{uw} +\sum_{uv}\langle uy||vw\rangle\gamma_{ux}\gamma_{vz} \nonumber\\
  &\quad -\sum_{uv}\langle uy||vz\rangle\gamma_{ux}\gamma_{vw} -\sum_{uv}\langle ux||vw\rangle\gamma_{uy}\gamma_{vz} +\sum_{uv}\langle ux||vz\rangle\gamma_{uy}\gamma_{vw} \nonumber\\
  &\quad -\sum_{uv}\langle uw||vy\rangle\lambda_{ux,vz} +\sum_{uv}\langle uz||vy\rangle\lambda_{ux,vw} +\sum_{uv}\langle uw||vx\rangle\lambda_{uy,vz} \nonumber\\
  &\quad -\sum_{uv}\langle uz||vx\rangle\lambda_{uy,vw}
  \label{eq:pp-aa-aa}
\end{align}

\paragraph*{Mixed blocks.}
\begin{align}
  M^{\pp}_{(ax),(by)} &= \delta_{xy}f_{ab} +\delta_{ab}f_{xy} -\delta_{ab}\bar{\mathcal L}_{xy} \nonumber\\
  &\quad -f_{ab}\gamma_{xy} +\langle ax||by\rangle +\sum_{u}\langle ua||by\rangle\gamma_{ux} \nonumber\\
  &\quad +\sum_{u}\langle ub||ax\rangle\gamma_{uy} +\sum_{uv}\langle ua||vb\rangle\gamma_{ux}\gamma_{vy} -\sum_{uv}\langle ub||va\rangle\lambda_{ux,vy}
  \label{eq:pp-av-av}\\
  M^{\pp}_{(ix),(jy)} &= -\delta_{ij}\bar{\mathcal L}_{xy} -f_{ij}\gamma_{xy} +\sum_{uv}\langle ui||vj\rangle\gamma_{ux}\gamma_{vy} \nonumber\\
  &\quad -\sum_{uv}\langle uj||vi\rangle\lambda_{ux,vy}
  \label{eq:pp-ca-ca}\\
  M^{\pp}_{(ab),(cx)} &= -\delta_{bc}f_{ax} +\delta_{ac}f_{bx} +\delta_{bc}\bar{\mathcal L}_{ax} \nonumber\\
  &\quad -\delta_{ac}\bar{\mathcal L}_{bx} +\langle ab||cx\rangle +\sum_{u}\langle uc||ab\rangle\gamma_{ux}
  \label{eq:pp-vv-av}\\
  M^{\pp}_{(ab),(xy)} &= \langle ab||xy\rangle -\sum_{u}\langle uy||ab\rangle\gamma_{ux} +\sum_{u}\langle ux||ab\rangle\gamma_{uy}
  \label{eq:pp-vv-aa}\\
  M^{\pp}_{(ab),(ix)} &= \sum_{u}\langle ui||ab\rangle\gamma_{ux}
  \label{eq:pp-vv-ca}\\
  M^{\pp}_{(ax),(yz)} &= \delta_{xz}f_{ay} -\delta_{xy}f_{az} -\delta_{xz}\bar{\mathcal L}_{ay} \nonumber\\
  &\quad +\delta_{xy}\bar{\mathcal L}_{az} -f_{ay}\gamma_{xz} +f_{az}\gamma_{xy} \nonumber\\
  &\quad +\langle ax||yz\rangle +\sum_{u}\langle ua||yz\rangle\gamma_{ux} -\sum_{u}\langle uz||ax\rangle\gamma_{uy} \nonumber\\
  &\quad +\sum_{u}\langle uy||ax\rangle\gamma_{uz} -\sum_{uv}\langle ua||vz\rangle\gamma_{ux}\gamma_{vy} +\sum_{uv}\langle ua||vy\rangle\gamma_{ux}\gamma_{vz} \nonumber\\
  &\quad +\sum_{uv}\langle uz||va\rangle\lambda_{ux,vy} -\sum_{uv}\langle uy||va\rangle\lambda_{ux,vz}
  \label{eq:pp-av-aa}\\
  M^{\pp}_{(ax),(iy)} &= \tfrac{1}{2}\delta_{xy}f_{ai} -f_{ai}\gamma_{xy} +\sum_{u}\langle ui||ax\rangle\gamma_{uy} \nonumber\\
  &\quad +\sum_{uv}\langle ua||vi\rangle\gamma_{ux}\gamma_{vy} -\sum_{uv}\langle ui||va\rangle\lambda_{ux,vy}
  \label{eq:pp-av-ca}\\
  M^{\pp}_{(ax),(ij)} &= -\langle ax||ij\rangle -\sum_{u}\langle ua||ij\rangle\gamma_{ux}
  \label{eq:pp-av-cc}\\
  M^{\pp}_{(xy),(iz)} &= \delta_{yz}f_{xi} -\delta_{xz}f_{yi} -\delta_{yz}\bar{\mathcal L}_{xi} \nonumber\\
  &\quad +\delta_{xz}\bar{\mathcal L}_{yi} -f_{xi}\gamma_{yz} +f_{yi}\gamma_{xz} \nonumber\\
  &\quad +\sum_{u}\langle ui||xy\rangle\gamma_{uz} -\sum_{uv}\langle uy||vi\rangle\gamma_{ux}\gamma_{vz} +\sum_{uv}\langle ux||vi\rangle\gamma_{uy}\gamma_{vz} \nonumber\\
  &\quad +\sum_{uv}\langle ui||vy\rangle\lambda_{ux,vz} -\sum_{uv}\langle ui||vx\rangle\lambda_{uy,vz}
  \label{eq:pp-aa-ca}\\
  M^{\pp}_{(xy),(ij)} &= -\langle xy||ij\rangle +\sum_{u}\langle uy||ij\rangle\gamma_{ux} -\sum_{u}\langle ux||ij\rangle\gamma_{uy}
  \label{eq:pp-aa-cc}\\
  M^{\pp}_{(ix),(jk)} &= \delta_{ik}\bar{\mathcal L}_{xj} -\delta_{ij}\bar{\mathcal L}_{xk} -\sum_{u}\langle ui||jk\rangle\gamma_{ux}
  \label{eq:pp-ca-cc}
\end{align}
Remarks.  (i) The ladder integral is dressed on each leg by $-\gamma$
[Eq.~\eqref{eq:pp-M}]; a core leg therefore subtracts one copy of the bare
integral, so that two core legs in a pair give
$\mathcal M^{\pp}_{(ab),(ij)}=-\langle ab||ij\rangle$ and one core leg gives the
$-\langle\cdot\cdot||i\cdot\rangle$ combinations, while virtual legs are
undressed.  (ii) The cumulant enters
only through the particle--hole-type contractions
$\langle u\cdot||v\cdot\rangle\lambda_{u\cdot,v\cdot}$ and through $\sLag$; it is
absent from all blocks whose ket or bra pair is purely inactive.  (iii)
The $\scv$ pairs are absent from the manifold but the $\sav$ and $\sca$
sectors couple the particle-like and hole-like spaces through the active
density: $\mathcal M^{\pp}_{(ax),(iy)}$ is the extended-vacuum analogue of a
$\bB^{\pp}$ block with one active index on each side.
\twocolumngrid

\subsection{Crossed-ring channel}
\label{sec:teh-blocks}

\Proposed By Eq.~\eqref{eq:teh-erpa} the crossed-ring channel has the same
sectors, the same metric blocks \eqref{eq:ph-S-blocks} and the same $A$ blocks
as the particle--hole channel, while every $B$ block changes sign,
$B^{\teh}=-B$, including the $B$ blocks of the $\saa$ sector obtained from
$\mathcal M_{(pq),(zw)}$ in the active natural-orbital basis.  Equivalently,
$\bcalM^{\teh}=\bZ\bcalM\bZ$ with $\bZ=+1$ on $(ai),(xi),(ax)$ and on the
active pairs with $n_y>n_x$, and $\bZ=-1$ on their partners; $\bZ$ acts within
each sector, so the sector structure is preserved.

%%%%%%%%%%%%%%%%%%%%%%%%%
\section{Numerical verification}
\label{sec:verification}
%%%%%%%%%%%%%%%%%%%%%%%%%

All equations of Secs.~\ref{sec:ph}--\ref{sec:blocks} were checked against
brute-force double commutators evaluated with explicit second-quantized
operators in the full Fock space (\texttt{numerics/erpa\_blocks/}):
\begin{enumerate}
\item[(a)] Eqs.~\eqref{eq:ph-S}--\eqref{eq:ph-asym} and
  \eqref{eq:pp-S}--\eqref{eq:pp-M-cum}, including the unsymmetrized forms,
  Eq.~\eqref{eq:brillouin}, the identity $\mathcal M_{(qp),(sr)}=\mathcal M_{(pq),(rs)}$
  and the cumulant form of $\Lag$, for random real $h$, $(pr|qs)$ with the
  eightfold symmetry and a random (non-eigenstate) three-electron vector in
  seven spin orbitals: maximum deviations $\le 2\times10^{-15}$
  (\texttt{verify\_general.py}).
\item[(b)] Eqs.~\eqref{eq:cas-rdm}--\eqref{eq:cas-lag}, \eqref{eq:ph-S-blocks},
  \eqref{eq:pp-S-blocks} and all 35 blocks of
  Secs.~\ref{sec:ph-blocks}--\ref{sec:pp-blocks}, for a product state with
  two core, four active (two electrons, random CI vector) and two virtual spin
  orbitals and random integrals, so that $\Lag$ is not symmetric: maximum
  deviations $\le 3\times10^{-15}$ (\texttt{blocks.py}, which also generates
  the \LaTeX{} of the blocks from Eqs.~\eqref{eq:ph-M-cum} and
  \eqref{eq:pp-M-cum} by substituting the index types).
\end{enumerate}
The Hartree--Fock limits quoted in Secs.~\ref{sec:ph}--\ref{sec:teh} are the
$\scv$--$\scv$, $\svv$--$\svv$, $\scc$--$\scc$ and $\svv$--$\scc$ blocks with
$\gamma^{\act}=0$.

%%%%%%%%%%%%%%%%%%%%%%%%
\bibliography{parquet-mr-hedin}
%%%%%%%%%%%%%%%%%%%%%%%%

\end{document}
